\section{Método demonstrado}
A demonstração utiliza inspeção estática de um caso único e seleção intencional de três registros já presentes na matriz de seis unidades; não houve nova execução do jogo, coleta com estudantes ou cálculo de frequência global. O código analisado pertence ao baseline identificado no registro de pesquisa, enquanto o protocolo e a matriz pertencem à documentação científica da PR 2.

O procedimento consiste em localizar o enunciado do GDD, definir o critério da unidade, identificar o símbolo correspondente no código e registrar uma classificação acompanhada de justificativa; revisão humana e evidência de execução permanecem campos separados. A ausência de parâmetro numérico é registrada como impossibilidade de verificar a fidelidade daquele parâmetro, e não como ausência da funcionalidade.

As citações são processadas por comandos autor-data de \texttt{natbib}, com grafia normal nas chamadas, em coerência com a atualização da norma de citações; os registros bibliográficos continuam sujeitos a conferência humana, especialmente perante a edição de referências de 2025 \citep{abnt2023,ufes2026}.

\subsection{Identificação do corpus}
O GDD foi fornecido pelo responsável pelo projeto como arquivo \emph{GDD -- Night Heist Escape Phase (Versão 0).pdf}, com identificação interna V2 / versão 0.2; a leitura foi textual, sem conferência de bytes ou paginação. O código inspecionado corresponde ao baseline \nolinkurl{72699f0c2752d029d8e8c3b48420c18a70dbe2a1}, e o quadro reutiliza os registros U01, U04 e U05 da matriz piloto. Essas informações identificam o material analisado e não são usadas como fundamentação externa.
